\section{結言}
この日の話題を貫くのは、モデル性能の拡大と、実務を実行するための仕組みの発展である。Hy4は視覚・コンピュータ操作を統合する方向を示し、Conductorの買収とMistralのデモは、複数エージェントの統制や専門作業の自動化に関心が集まっていることを示した。

Claudeを使った構造生物学研究は、エージェント技術の応用先がソフトウェア開発にとどまらないことを印象づける。一方、モデル公開予定、実験の再現性、実運用での信頼性など、投稿だけでは結論を出せない点も残る。次の注目点は、オープンウェイトの実際の公開と、独立した評価による性能・安全性の確認だろう。

なお、OpenAI内部モデルに関する未確認の噂も観測されたが、Grokは公式確認がないとして主要トピックに採用していない。本号はX上の観測記録の編纂であり、独立した一次資料検証は行っていない。
